

%%% TIKZ parameters and macros

\newcommand{\rootShape}{ellipse}
\newcommand{\rootFill}{blue!15}
\newcommand{\elShape}{ellipse split}
\newcommand{\elFill}{green!15}
\newcommand{\txtShape}{ellipse split}
\newcommand{\txtFill}{red!10}
\newcommand{\attrShape}{ellipse split}
\newcommand{\attrFill}{red!20}

\newcommand{\textWidthTikz}{1cm}
\newcommand{\tikzLevelDistance}{15mm}
\newcommand{\tikzSiblingDistanceOne}{50mm}
\newcommand{\tikzSiblingDistanceTwo}{35mm}
\newcommand{\tikzSiblingDistanceThree}{15mm}
\newcommand{\tikzSiblingDistanceFour}{15mm}

%%
%% @par1 = id @par2=name
\newcommand{\tikzElementWithID}[2]{node[\elShape,fill=\elFill] (#1)
{\tikzEl{#2}}}
\newcommand{\tikzElement}[1]{node[\elShape,fill=\elFill] {\tikzEl{#1}}}

\newcommand{\tikzText}[1]{node[\txtShape,fill=\txtFill] {\tikzTx{#1}}}
\newcommand{\tikzAttr}[1]{node[\attrShape,fill=\attrFill,dotted] {\tikzAt{#1}}}

\newcommand{\tikzRoot}{\domtype{Document}}


\newcommand{\tikzEl}[1]{%
  \domtype{Element}
   \nodepart{lower} {#1}
}

\newcommand{\tikzTx}[1]{%
  \domtype{Text}
   \nodepart{lower} {#1}
}

\newcommand{\tikzAt}[1]{%
  \domtype{Attr.}
   \nodepart{lower} {#1}
}

%
% A DOM tree in tikz, in a floating figure
%
\newcommand{\tikzDomTree}[2]{%
\begin{figure}
\begin{center}

\begin{small}
\input{../textrees/#1.tikz}
\end{small}
   \caption{\label{#1} #2}
   
\end{center}
  \end{figure}
}

%
% Without the floating env.
%
\newcommand{\tikzDomTreeHere}[1]{%
\begin{center}

\input{../textrees/#1.tikz}

\end{center}
}


%%
%% Style definition
%

\tikzstyle{every text node part}=[font=\bf]
\tikzstyle{every lower node part}=[font=\itshape]

%% Tikz style for drawing DOM trees
\tikzstyle{domtree}=[fill=red!60,rectangle,inner sep=1pt,level
distance=\tikzLevelDistance]

%% Tikz style for ontologies
\tikzstyle{onto}=[inner sep=1pt,level
distance=\tikzLevelDistance, sibling distance=15mm]

%% ikz style for the root node
\tikzstyle{rootnode}=[\rootShape,fill=\rootFill]

 \tikzstyle{level 1}=[sibling distance=\tikzSiblingDistanceOne]
  \tikzstyle{level 2}=[sibling distance=\tikzSiblingDistanceTwo]
  \tikzstyle{level 3}=[sibling distance=\tikzSiblingDistanceThree]
  \tikzstyle{level 4}=[sibling distance=\tikzSiblingDistanceFour]
  